\documentclass[a4paper,11pt]{article}
\usepackage{jheppub}
\usepackage{amsmath}
\usepackage{amssymb}
\usepackage{lineno}
\linenumbers

\title{\boldmath Hidden in Plain Sight Since 1928:\\
The Gravitational Coupling Polynomial
$f(K) = K^0 + K^1 + K^2$}

\author{Stephen Nelson}
\affiliation{Field Dynamics Research\\
\href{mailto:contact@fielddynamics.org}{contact@fielddynamics.org}}

\preprint{Zenodo: [DOI to be inserted]}

\abstract{
The tetrad equation of general relativity,
$g_{\mu\nu} = \eta_{ab}\,e^a_\mu\,e^b_\nu$, has been in
continuous use since Einstein wrote it in 1928. We show
that it contains an unread polynomial. A rank-level scan
of the equation's structural content on the local frame
space identifies three occupied tensor ranks --- a
closure condition (rank~0), a frame (rank~1), and a
coupling template (rank~2) --- whose rank-dimensions sum
to $K^0 + K^1 + K^2 = 1 + 4 + 16 = 21$, where $K = 4$ is
the frame dimension. Two independent arguments then
establish that $K = 4$ is the unique value compatible with
field closure in $D = 3$ spatial dimensions: the combinatorial
self-duality of the tetrahedron via Pascal's triangle, and
three independent physical constraints from gravitational
wave observations, the simplex theorem, and ghost-free
bimetric gravity. The resulting polynomial
$f(K,s) = 1 + sK + K^2$ evaluates to three integers ---
21, 17, and 13 --- in its three structural states, locked
by an algebraic identity that permits no adjustment. A
companion letter develops the application of the
polynomial's contact coupling at $s = 0$ to the underdetermined
parameter space of Hassan--Rosen bimetric gravity.
}

\begin{document}
\maketitle

%--------------------------------------------------------------
\section{An Equation Used for a Century Without Being Read}
%--------------------------------------------------------------

Tetrad formalism is one of the oldest and most widely used
languages in gravitational physics. Einstein introduced
the tetrad in 1928~\cite{Einstein1928} as a way to work
with local inertial frames in curved spacetime. Since then,
it has become indispensable across several distinct
traditions. Regge calculus discretizes spacetime using the
tetrahedron as its fundamental three-dimensional building
block~\cite{Regge1961}; spinfoam models use the tetrahedral
intertwiner as the fundamental vertex of quantum
geometry~\cite{Rovelli}; and bimetric gravity employs two
independent vierbein fields $e^a_\mu$ and $\tilde{e}^a_\mu$
coupled through a shared interaction
structure~\cite{HR2012a,Alexandrov2013}. These frameworks
draw on different aspects of the same geometric source: the
local frame, the tetrahedral symmetry of three-dimensional
space, and the geometric condition that holds the structure
complete at every point. The tetrahedron and the tetrad are not peripheral
tools. They are at the core of how modern physics describes
the geometry of space.

A requirement common to all of these frameworks is
\emph{field closure}: the condition that the geometric
structure at each point is complete, non-degenerate, and
self-consistent --- that the field holds together rather
than collapsing. In Regge calculus, field closure means a
simplicial complex in which every boundary face is shared
by exactly two simplices: no open edges, no dangling
geometry, the structure complete~\cite{Regge1961}. In loop
quantum gravity and spinfoam models, it is the closure
constraint $\sum_i \vec{L}_i = 0$ at each spin network
vertex --- the Gauss law of the discrete geometry, ensuring
that the dual polyhedron closes~\cite{Rovelli}. In tetrad
gravity, it is $\det(e) \neq 0$: the legs of the local
frame span the tangent space completely, and the metric is
non-degenerate. In every case the logic is the same: the
field either closes --- the geometry holds --- or it does
not. Field closure is not a condition imposed from outside.
It is the minimal requirement that the geometry be real.

In the tetrad formulation of bimetric gravity, both
vierbein fields $e^a_\mu$ and $\tilde{e}^a_\mu$ must be
simultaneously defined at each spacetime point, and the
ghost-free interaction potential is built from the
transition matrix $S^a{}_b = (e^{-1}\tilde{e})^a{}_b$
evaluated at that shared point~\cite{HR2012a,Alexandrov2013}.
Standard bimetric treatments leave the relative spatial
orientation of the two vierbein fields as an unconstrained
choice. The closure condition derived in Section~3 removes
this freedom: non-degenerate field closure requires the two
sectors to carry opposite spatial orientations,
$\mathrm{sgn}(\det e) = -\mathrm{sgn}(\det \tilde{e})$,
so that each defines its own tetrahedral frame at the
coupling point. We call the coupling state at which
the two sectors meet --- the $s = 0$ evaluation of the
polynomial introduced in Section~3 --- the \emph{Field
Origin}. We call the sector containing the physical
observer the \emph{primary sector}; its geometrically
required partner with reversed spatial orientation is
the \emph{dual sector}. The resulting geometry is two
interlocked tetrahedral sectors of opposite orientation,
connected through their shared Field Origin.

At the foundation of all of this sits a single equation:
\begin{equation}
  g_{\mu\nu} = \eta_{ab}\,e^a_\mu\,e^b_\nu.
  \label{eq:tetrad}
\end{equation}
It appears in every standard textbook on gravitational
physics~\cite{MTW,Carroll,Rovelli} and has been treated,
for nearly a century, as a definition: a convenient way to
encode the spacetime metric through a locally defined
orthonormal frame. The internal metric $\eta_{ab}$ is the
Minkowski metric of special relativity:
$\eta_{ab} = \mathrm{diag}(-1,+1,+1,+1)$. This is the
Lorentzian signature of modern physics --- one timelike
direction and three spacelike directions. It is not an
assumption of the tetrad formalism. It is the signature of
spacetime itself, the geometric fact from which the
$3+1$ structure of general relativity descends. The tetrad
$e^a_\mu$ assigns one orthonormal leg to each of these
four directions: one timelike leg $e^0_\mu$ and three
spacelike legs $e^i_\mu$, $i = 1,2,3$. The frame dimension
is therefore $K = 4$ --- fixed not by choice but by the
Lorentzian geometry of spacetime.

Physicists use it, differentiate it, expand perturbation
theory around it, and derive consequences from it. They have
been very good at answering one particular question about it:
how many independent components does $g_{\mu\nu}$ have?
The answer is ten.

That is the right answer to that question. But it is not
the question this paper asks.

The question here is different, and simpler: how many
structurally distinct \emph{rank levels} does
Eq.~\eqref{eq:tetrad} occupy on the local frame space?
Not how many numbers fill each slot, but how many tensor
ranks carry non-trivial content --- and what is the
rank-dimension at each occupied level? When you ask this
question, as we explain in the next section, the answer
is $1 + 4 + 16 = 21$.

This sum is the coupling polynomial $f(K) = 1 + K + K^2$
evaluated at $K = 4$. It has been present in
Eq.~\eqref{eq:tetrad} since 1928. It has not been read
before because the standard question in tetrad gravity
asks about components, not structural types. This paper
reads it --- and then shows, through two independent
arguments, that no other polynomial could have been
hiding there.


%--------------------------------------------------------------
\section{Reading the Polynomial: Rank-Dimension Counting}
%--------------------------------------------------------------

Before we can count, we need to agree on what we are
counting. There are three natural ways to measure the
``size'' of an object in the tetrad equation, and they
give three different answers. Choosing the wrong one
produces a number that looks correct but carries no
structural meaning. Choosing the right one produces the
polynomial.

The first measure is \emph{independent components}: after
imposing all symmetries, how many free numbers characterize
the object? For the symmetric metric $g_{\mu\nu}$, the
answer is ten. For the tetrad $e^a_\mu$, the answer is
sixteen, reduced to ten by local Lorentz freedom --- the
freedom to boost and rotate the orthonormal frame
independently at each spacetime point. These are correct answers to a question about
physical degrees of freedom.

The second measure is \emph{nonzero entries}: for the flat
Minkowski metric $\eta_{ab} = \mathrm{diag}(-1,+1,+1,+1)$,
the matrix has $4 \times 4 = 16$ entries in total, of which
four are nonzero diagonal entries and twelve are zeros.
This correctly describes a particular signature choice.
It counts how the sixteen available pairings happen to be
weighted in flat spacetime.

The third measure is \emph{rank-dimension}: for a tensor
of rank $r$ defined on the local frame space --- the
four-dimensional space indexed by the frame index $a$,
of dimension $K$ --- the rank-dimension is simply $K^r$. It does not depend on
which entries are zero, which are equal, or which symmetry
has been imposed. It counts how many frame-direction
combinations the object must be defined on --- the size
of the template itself, before any values are assigned.

The polynomial asks the third question. More precisely,
it performs a \emph{rank-level scan}: at each tensor rank
$r = 0, 1, 2, \ldots$ on the local frame space, it asks
whether the tetrad equation requires a non-trivial
structural condition at that rank. If a rank is occupied
--- if the theory demands something there --- it
contributes $K^r$ to the sum. If a rank is unoccupied,
it contributes nothing. We now scan the ranks.

\textbf{Rank~0: The closure condition.}
A rank-0 object on frame space is a scalar: a single
number with no free frame index. In the vierbein
formulation of gravity, the action takes the form
$S = \int e\,\mathcal{L}\,d^4x$, where
$e = \det(e^a_\mu)$ is the scalar volume measure --- the
Jacobian that converts coordinate volume into geometric
volume~\cite{Carroll,MTW}. The requirement $e \neq 0$
is the statement that the rank-0 sector of the frame
bundle is non-trivially occupied: the geometry closes
rather than degenerates.

This point deserves emphasis, because a reader may
object that $\det(e)$ is computed from $e^a_\mu$ and
therefore not a separate structural ingredient. But
consider a soap bubble. The closure of the film ---
the condition that it forms a bounded surface enclosing
volume --- is not a separate material added to the soap.
It is produced by the film. Yet a soap film that does
not close is a fundamentally different object: an open
surface, not a bubble. Closure is structurally
independent of the film's local properties in the
precise sense that you can specify the film everywhere
and still fail to have closure. The same is true here.
One can write down a tetrad $e^a_\mu$ with all sixteen
components specified and still have $\det(e) = 0$ at a
point. At that point the frame does not span the tangent
space, the metric degenerates, and the equation
$g_{\mu\nu} = \eta_{ab}\,e^a_\mu\,e^b_\nu$ describes
no geometry. Moreover, $\det(e) \neq 0$ is not external to
Eq.~\eqref{eq:tetrad}: the left-hand side $g_{\mu\nu}$ is
a non-degenerate metric only when $\det(e) \neq 0$. The
rank-0 condition is the precondition for the equation to
mean what it says. It is the independent structural
requirement that the geometry holds --- that the field
closes rather than collapses. It contributes one
structural element to the count.
\begin{equation*}
  \text{Rank-dimension of rank 0:}\quad K^0 = 1.
\end{equation*}

\textbf{Rank~1: The frame.}
A rank-1 object on frame space is a vector: one free
frame index $a$, running over $K$ values. The tetrad
$e^a_\mu$ carries exactly one free frame index. It assigns
a spacetime one-form to each of the $K$ frame
directions: one timelike and three spacelike. This is
the frame --- the set of directions the local geometry
provides. It contributes $K$ structural elements.
\begin{equation*}
  \text{Rank-dimension of rank 1:}\quad K^1 = K = 4.
\end{equation*}

\textbf{Rank~2: The coupling template.}
A rank-2 object on frame space is a matrix: two free
frame indices $a$ and $b$, each running over $K$ values.
The contraction $\eta_{ab}\,e^a_\mu\,e^b_\nu$ evaluates
one term for every ordered pair $(a,b)$. There are
$K \times K$ such pairs.

This step deserves care, because it is where the standard
component count and the rank-dimension count diverge.
The flat Minkowski metric $\eta_{ab}$ assigns $\pm 1$
to the four diagonal pairs and zero to the twelve
off-diagonal pairs. It is tempting to read those zeros
as absences --- as if only four pairings are really
active. But that reading is wrong. The zeros are not
absences. They are assignments of zero weight to twelve
of the sixteen pairings. The template itself --- the
double contraction over all ordered pairs $(a,b)$ ---
evaluates all sixteen regardless of what values
$\eta_{ab}$ takes. A Euclidean metric would assign $+1$
across the same sixteen pairings. A different signature choice would distribute those sixteen
values differently still. What is structurally
invariant is the template: $K \times K$ frame-direction
pairings that the coupling must be defined on.
\begin{equation*}
  \text{Rank-dimension of rank 2:}\quad K^2 = 16.
\end{equation*}

Summing the occupied ranks by rank-dimension:
\begin{equation}
  \underbrace{K^0}_{= 1,\ \text{existence}}
  +\underbrace{K^1}_{= 4,\ \text{frame}}
  +\underbrace{K^2}_{= 16,\ \text{pairing template}}
  = 1 + 4 + 16 = 21.
  \label{eq:sum}
\end{equation}

This sum is the total dimension of the graded frame algebra
$\bigoplus_{r=0}^{2} V^{(r)}$, where $V^{(r)}$ is the
space of rank-$r$ tensors on the $K$-dimensional frame
space --- the same operation that yields the Euler
characteristic when the alternating-sign convention is
applied to a simplicial complex~\cite{Regge1961}. Summing
dimensions across tensor ranks is not exotic; it is one of
the most standard operations in algebra and topology.

The coupling polynomial $f(K) = K^0 + K^1 + K^2$ is this
graded dimension evaluated at $K = 4$. The polynomial counts how many
structural levels the field must traverse before field
closure is achieved --- before the geometry holds completely.
We call this the \emph{coupling depth}; at $K = 4$ in
three spatial dimensions, the coupling depth of the
tetrad equation is 21.

The scan terminates at rank~2. The tetrad equation has
exactly three occupied ranks because a tetrad theory
requires exactly three levels of structural description:
whether spacetime is non-degenerate (rank~0), which
directions exist (rank~1), and how those directions
couple to each other (rank~2). This mirrors the standard
decomposition in gauge field theory --- scalar sector,
vector sector, tensor sector --- and for the same
reason: these are the only ranks at which the field
carries independent structural content. A rank-3 object
would require three simultaneous free frame indices at
each coupling event. No such object appears in
$g_{\mu\nu} = \eta_{ab}\,e^a_\mu\,e^b_\nu$, in the
tetrad action, or in the field equations of any metric
or bimetric gravity theory. The equation closes at
rank~2, and the polynomial closes with it.

One clarifying note. When independent components are
counted instead of rank-dimensions --- one for the
scalar volume measure, ten for the tetrad after fixing
local Lorentz freedom, ten for the symmetric metric --- the sum is
$1 + 10 + 10 = 21$: the same number, by coincidence.
That coincidence is worth naming because a reader might
wonder whether both counts say the same thing. They
do not. The independent-component count reaches 21 by
accident, with no structural reason for the total. Only
the rank-dimension count is a polynomial in $K$; the
component count $1 + 10 + 10$ does not generalize to
other values of $K$ and carries no algebraic structure.
The polynomial emerges from the rank-dimension count.
It does not emerge from the component count.

The polynomial $f(K) = K^0 + K^1 + K^2$ was written in
1928 as $g_{\mu\nu} = \eta_{ab}\,e^a_\mu\,e^b_\nu$.
It has taken a century to read it.


%--------------------------------------------------------------
\section{Why $K = 4$: The Closure Condition}
%--------------------------------------------------------------

Section~2 presented the polynomial form with $K$ as a free
symbol to keep the structural argument general. This section
establishes $K = 4$ by a second, independent route: the
combinatorial self-duality of the tetrahedron and the
closure condition it imposes. The specific polynomial with
values $1 + 4 + 16 = 21$ requires this.

The tetrahedron is the three-dimensional simplex: the
minimal solid that encloses a bounded region of space
with the fewest possible faces. Its complete anatomy ---
every vertex, edge, face, and body --- is encoded in
row $K = 4$ of Pascal's triangle:
\begin{equation}
  \binom{4}{0} = 1,\quad
  \binom{4}{1} = 4,\quad
  \binom{4}{2} = 6,\quad
  \binom{4}{3} = 4,\quad
  \binom{4}{4} = 1.
\end{equation}
These five numbers count, in order: the empty set
(the existence condition), vertices, edges, faces, and
the enclosed body. They are determined entirely by $K = 4$
and carry no free choices.

Look at the row carefully: 1, 4, 6, 4, 1. It reads the
same forwards and backwards. The second and fourth entries
are equal --- both are 4. In geometric terms, the
tetrahedron has the same number of vertices as it has
faces. This is \emph{self-duality}. The dual of a
polyhedron is built by placing a new vertex at the center
of each face; for the tetrahedron, the result is another
tetrahedron. Vertices and faces are structurally
interchangeable. No other regular solid has this property.

Self-duality forces the polynomial into exactly three
terms. In Regge calculus, each Pascal entry corresponds
to a distinct geometric stratum of the simplex ---
vertices, edges, faces, body --- each contributing an
independent coupling channel to the interaction
structure~\cite{Regge1961}. If vertex and face counts
were different --- as they are for the cube, the
octahedron, or any other regular polyhedron --- a coupling
formula built from Pascal's row would need separate terms
for vertex channels and face channels, producing five
terms in all. Because the tetrahedron is self-dual, those
separate terms collapse into a single propagation term.
Five Pascal entries become three polynomial terms:
\begin{equation}
  K^0 = 1 \quad (\text{existence}),\qquad
  K^1 = 4 \quad (\text{propagation}),\qquad
  K^2 = 16 \quad (\text{pair-complete interaction: every
  ordered pair $(a,b)$ of frame directions contributing
  one coupling channel}).
\end{equation}

Now we apply the closure condition. The two tetrad sectors
established in Section~1 --- each carrying a tetrahedral
frame, meeting at the Field Origin --- form the geometric
object whose closure we must enforce. Mapping from the
primary sector to the dual sector is a spatial orientation
reversal: $\mathrm{sgn}(\det e)$ flips sign. We introduce
a state parameter $s \in \{+1, 0, -1\}$ to index this
orientation. To evaluate the polynomial in the dual sector,
$s$ must act on each term according to the transformation
rule for its tensor rank under spatial reflection.

Section~2 established three occupied ranks.
Their behavior under spatial reflection follows directly
from that classification:
\begin{itemize}
\item \textbf{Rank-0 (scalar volume measure, $K^0 = 1$):}
A scalar has no spatial extent or direction. It is strictly
invariant under spatial orientation reversal. The term
contributes $+1$ in both sectors.
\item \textbf{Rank-1 (tetrad frame, $K^1 = K$):}
A vector carries one free frame index and therefore
possesses directionality. Under spatial orientation
reversal, the spatial components of a vector flip sign.
The term contributes $+K$ in the primary sector and
$-K$ in the dual sector: it carries the factor $s$.
\item \textbf{Rank-2 (coupling template, $K^2$):}
A rank-2 tensor carries two free frame indices. Under
spatial orientation reversal, each index contributes a
sign flip, and $(-s)(-s) = +1$. The bilinear coupling
template is even under parity and therefore invariant.
The term contributes $+K^2$ in both sectors.
\end{itemize}
The state parameter $s$ therefore enters the polynomial
exclusively in the linear term, by geometric necessity:
\begin{equation}
  f(K, s) = 1 + sK + K^2.
  \label{eq:polynomial_s}
\end{equation}
This is not a parametric choice. It is the unique form
consistent with how rank-0, rank-1, and rank-2 objects
transform under spatial reflection. This classification
follows the standard representation-theoretic convention
for spatial inversion in Lorentzian spacetime: a rank-1
object has parity $-1$ as a representation, just as a
four-vector in quantum field theory is classified as
parity-odd even though its timelike component is invariant
under spatial reflection. The Lorentzian signature
$(-1,+1,+1,+1)$ is precisely what makes the $3+1$ split
--- and therefore the definition of spatial parity ---
well-defined. The invariant timelike direction is the
anchor that gives spatial orientation its meaning, not a
counterexample to the parity assignment.

For this dual-sector structure to be topologically
closed, the dynamic range of the coupling polynomial
across its two chambers must match the combinatorial
content of the simplex boundary. This requirement has a
concrete algebraic counterpart. The Hassan--Rosen
ghost-free interaction potential is built from the
elementary symmetric polynomials $e_n(S)$ of the frame
transition matrix
$S^a{}_b$~\cite{HR2012a,HR2012b,Alexandrov2013}. In the
canonical (Hamiltonian) formulation, the physically
independent interaction structure is carried by the $D$
spatial eigenvalues of $S$, whose elementary symmetric
polynomials organize the $2^D$ subsets of the spatial
simplex boundary by
cardinality~\cite{Alexandrov2013}. For the dynamic range
of the coupling polynomial to encode all and only these
subsets --- one coupling channel per boundary subset,
with no degeneracy and no excess --- requires
$\Delta f = 2^D$. If $\Delta f < 2^D$, boundary subsets
share coupling channels and information is lost at the
boundary: the closure is degenerate. If
$\Delta f > 2^D$, coupling channels exist with no
boundary subset to absorb them --- the potential
contains interaction terms beyond those permitted by the
elementary symmetric polynomial structure, which is
precisely the condition that reintroduces the
Boulware--Deser ghost~\cite{BD1972,HR2012b}. The unique
condition for non-degenerate, ghost-free closure is
therefore $\Delta f = 2^D$.

The dynamic range of the polynomial across the two tetrad
sectors is the difference between the primary and dual
evaluations:
\begin{equation}
  \Delta f = f(K,+1) - f(K,-1) = (1 + K + K^2) - (1 - K + K^2) = 2K.
  \label{eq:range}
\end{equation}
This follows from the polynomial form alone --- no
assumption about the value of $K$. Note that the rank-0
and rank-2 contributions cancel identically in the
difference. This is not an accident: the parity
classification established above guarantees that only the
parity-odd sector (rank~1) distinguishes one tetrad
sector from the other. The parity-even terms carry no
sector-distinguishing information and therefore contribute
nothing to the dynamic range. The full polynomial is
nevertheless essential: it is the rank-level scan that
establishes $K$ as the rank-dimension of the unique
parity-odd level, and it is the parity structure of the
complete polynomial that isolates $\Delta f = 2K$ as the
dynamic range. Setting the dynamic
range equal to the boundary combinatorial content:
\begin{equation}
  \Delta f = 2^D
  \quad\Longrightarrow\quad
  2K = 2^D.
  \label{eq:closure}
\end{equation}
Substituting the simplex relation $K = D + 1$ ---
the vertex count of the $D$-simplex, which equals the
frame dimension of the tetrad in $D$ spatial dimensions
(Confirmation~2 below establishes this identification
independently):
\begin{equation}
  2(D + 1) = 2^D.
  \label{eq:diophantine}
\end{equation}
This is a Diophantine equation in $D$. Checking all
positive integers exhaustively:
$D = 1$: $4 \neq 2$;\quad
$D = 2$: $6 \neq 4$;\quad
$D = 3$: $8 = 8$ \checkmark;\quad
$D = 4$: $10 \neq 16$;\quad
and for all $D \geq 4$ the right side grows exponentially
while the left grows linearly, so no further solution
exists. The equation has exactly one positive integer
solution: $D = 3$, and therefore $K = 4$.

$D = 3$ is not a physical postulate. It is the unique
integer dimension at which the dual tetrad sector geometry
closes. This single result fixes both the spatial simplex
face count and the spacetime frame dimension simultaneously:
the $D$-simplex has $K = D+1 = 4$ faces, and spacetime has
$D+1 = 4$ dimensions --- the same $K = 4$ from the same
$D = 3$, not two separate assumptions. With $K = 4$ fixed
by this condition, the polynomial is fully determined:
$f(K) = 1 + 4 + 16 = 21$. The form came from the tetrad
equation in Section~2. The value of $K$ comes from the
closure of the dual-sector geometry here. Together they
give the complete polynomial with no freedom remaining.


%--------------------------------------------------------------
\section{Why $K = 4$: Three Physical Confirmations}
%--------------------------------------------------------------

The closure condition fixes $K = 4$ from pure geometry.
Three independent physical facts --- each from a different
domain of physics --- confirm it from the external world.
Their agreement is not required by the geometry, which
makes it a meaningful check: the universe has chosen
the same value that the tetrahedron requires.

\paragraph{Confirmation~1: Gravitational wave polarizations.}
The LIGO--Virgo detection of GW170817~\cite{LIGO2017}
established that exactly two tensor polarization modes
propagate through spacetime. The number of propagating
gravitational degrees of freedom is obtained by
subtracting $K$ Hamiltonian constraints and $K$
diffeomorphism constraints from the ten independent
components of the metric. This subtraction is not
coincidental: the $K$ basis vectors of the tetrad
directly generate $K$ diffeomorphism freedoms through
spacetime translations, and $K$ first-class Hamiltonian
constraints through the local Lorentz symmetry of the
frame~\cite{MTW,Carroll}. The result is:
\begin{equation}
  \text{propagating d.o.f.} = 10 - 2K.
\end{equation}
Setting this equal to 2 gives $K = 4$ uniquely. The
observation that gravitational waves carry exactly two
polarizations is equivalent to the statement that the
local frame has exactly four legs.

\paragraph{Confirmation~2: The simplex theorem.}
In $D$ spatial dimensions, the minimum number of points
needed to enclose a bounded region of space is $D + 1$~\cite{MTW}.
This is a theorem, not an assumption. In three spatial
dimensions, the minimum enclosing solid therefore has
$K = D + 1 = 4$ vertices and, by self-duality, $K = 4$
faces. Any coupling structure that operates in
three-dimensional space must have $K = 4$ channels.
This and Confirmation~1 reach $K = 4$ from entirely
unrelated starting points; their agreement is a
non-trivial consistency check.

\paragraph{Confirmation~3: Ghost-free bimetric coupling.}
Boulware and Deser showed~\cite{BD1972} that theories
coupling two interacting spin-2 fields generically
propagate an unphysical mode with negative kinetic energy,
destabilizing the vacuum. Hassan and Rosen
proved~\cite{HR2012a,HR2012b} that this pathology is
eliminated if and only if the interaction potential is
built from elementary symmetric polynomials of the frame
transition matrix between two tetrad sectors. This
independently establishes the dual-tetrad geometry ---
two interacting $K = 4$ frame fields coupled through a
shared transition matrix --- as the necessary architecture
for ghost-free spin-2 coupling. Standard bimetric gravity
leaves the relative orientation of the two sectors
unconstrained; the closure condition derived in Section~3
selects the opposite-orientation branch as the unique
non-degenerate configuration.

\paragraph{From confirmations to polynomial.}
Each confirmation independently requires $K = 4$; together
they confirm the complete polynomial structure. The
existence condition ($K^0 = 1$) is the requirement that
the coupling geometry be non-degenerate --- that the
field close at all. Ghost-free bimetric gravity
independently confirms that a coupling structure between
two tetrad sectors must exist. The propagation term
($K^1 = 4$) is the face count fixed by the simplex
theorem and confirmed by LIGO's two-polarization
measurement. The pair-complete interaction term
($K^2 = 16$) is the complete frame transition: every
face of one tetrahedral sector paired with every face
of the other, $4 \times 4 = 16$ independent coupling
directions --- the geometric realisation of which, at the
dual tetrahedral boundary, is developed in the companion
letter~\cite{Nelson2026b}. No term can be added, because the ghost-free
interaction potential acts on the tetrahedral boundary
--- the two-dimensional surface of faces --- as
established in the Hassan--Rosen canonical
analysis~\cite{Alexandrov2013}. A $K^3$ coupling would
require a three-dimensional cell on that two-dimensional
boundary, which does not exist. No term can be removed,
because each one is demanded by a separate physical
requirement. The polynomial $1 + sK + K^2$ with $K = 4$
satisfies all three confirmations simultaneously.


%--------------------------------------------------------------
\section{The Arithmetic Mean Lock}
%--------------------------------------------------------------

The polynomial is now fully established: $f(K,s) = 1 + sK + K^2$
with $K = 4$. A natural next question is how rigid this
result is. Could the integers it produces be adjusted, even
slightly, without breaking something?

The answer is no. The two interlocked tetrahedral sectors
established in Section~1 --- each carrying opposite spatial
orientation, meeting at their shared Field Origin --- give
the polynomial three natural evaluation points. The state
parameter $s$ established in Section~3 takes three values:
$s = +1$ in the primary sector, where the observer sits
and the full gravitational response is active; $s = 0$ at
the Field Origin, where both sectors meet at a single
shared point, spatial extent collapses to zero, and the
propagation term contributes nothing; and $s = -1$ in the
dual sector, where the spatial orientation is reversed.
The three evaluations are:
\begin{equation}
  f(4,+1) = 21, \qquad f(4,0) = 17, \qquad f(4,-1) = 13.
  \label{eq:states}
\end{equation}

These three integers are not independent. For any polynomial
of the form $1 + sK + K^2$, evaluating at $s = +1$, $0$,
and $-1$ always produces three equally spaced values, with
spacing exactly $K$. The middle value is always the
arithmetic mean of the outer two:
\begin{equation}
  \frac{f(K,+1) + f(K,-1)}{2} = f(K,0)
  \quad\Longrightarrow\quad
  \frac{21 + 13}{2} = 17.
  \label{eq:lock}
\end{equation}
This is not a coincidence specific to $K = 4$. It is an
algebraic consequence of the polynomial's form, forced
for any $K$.

Think of the three integers as equally spaced marks on a
ruler, with spacing $K = 4$. Moving any one of them
requires moving all three, because they are not three
separate numbers --- they are one polynomial evaluated
at three points. More precisely: if the Field Origin
value were 18 instead of 17, the equal spacing would
break, the primary state $1 + K + K^2$ would cease to
be the finite geometric series $\sum_{i=0}^{D-1} K^i$ for
$D = 3$, causing the geometric series identity
\begin{equation}
  f(K) = \frac{K^D - 1}{K - 1}
\end{equation}
to fail, and the connection between the polynomial
and the dimension of space $D = 3$ would be severed.
The integers 21, 17, and 13 form a single rigid object.
There are no free parameters to adjust.


%--------------------------------------------------------------
\section{Discussion and Conclusion}
%--------------------------------------------------------------

The tetrad equation $g_{\mu\nu} = \eta_{ab}\,e^a_\mu\,e^b_\nu$
is one of the most-used equations in gravitational physics.
Regge calculus builds discrete spacetime from tetrahedra.
Bimetric gravity couples two vierbein fields through a dual
tetrad structure~\cite{Alexandrov2013}. Spinfoam models use
the tetrahedral intertwiner as the building block of quantum
geometry. Every one of these formalisms has been working with
the ingredients of the coupling polynomial. None had asked
what those ingredients sum to when counted by rank-dimension.

The answer is $K^0 + K^1 + K^2 = 1 + 4 + 16 = 21$: the
coupling polynomial $f(K) = 1 + K + K^2$, uniquely fixed
by the Diophantine closure condition $2(D+1) = 2^D$ at
$D = 3$, $K = 4$, and confirmed by three independent
physical facts. The coupling
depth of the tetrad equation is 21. The three
structural integers 21, 17, and 13 are locked by the
arithmetic mean identity $(21 + 13)/2 = 17$ and form a
single rigid object with no adjustable parameters.

The polynomial evaluated at the Field Origin under contact
coupling, $s = 0$, gives $f(K,0) = 1 + K^2 = 17$, whose
interaction term $K^2 = 16$ constrains the weight of the
second-order coupling in the Hassan--Rosen bimetric
interaction potential. A companion letter~\cite{Nelson2026b}
shows that this evaluation, combined with orientation parity
and exchange symmetry, uniquely locks all five parameters
of the Hassan--Rosen bimetric interaction potential to
$\{\beta_n\} = \{1, 0, 16, 0, 1\}$, producing a testable
prediction at the MOND transition scale.

The coupling polynomial $f(K,s) = 1 + sK + K^2$ is not
new. It is among the oldest objects in gravitational
physics, encoded in the tetrad equation since 1928.
Tetrad formalism, Regge calculus, spinfoam theory, and
bimetric gravity have all been using its ingredients
without reading the sum. The question was always there.
This paper asks it.


%--------------------------------------------------------------
\acknowledgments
%--------------------------------------------------------------

The author acknowledges the tetrad gravity formalism of
Cartan, Einstein, and Weyl; the simplicial framework of
Regge; and the ghost-free bimetric structure of Hassan
and Rosen. Thanks to the broader community working on
the geometric foundations of gravitational field theory.


%--------------------------------------------------------------
\begin{thebibliography}{99}
%--------------------------------------------------------------

\bibitem{Einstein1928}
A.~Einstein,
\emph{Riemann-Geometrie mit Aufrechterhaltung des Begriffes
des Fernparallelismus},
Sitzungsber.\ Preuss.\ Akad.\ Wiss.\ (1928) 217.

\bibitem{Regge1961}
T.~Regge,
\emph{General relativity without coordinates},
Nuovo Cimento \textbf{19} (1961) 558.

\bibitem{MTW}
C.~W.~Misner, K.~S.~Thorne and J.~A.~Wheeler,
\emph{Gravitation},
W.~H.~Freeman, San Francisco (1973).

\bibitem{Carroll}
S.~M.~Carroll,
\emph{Spacetime and Geometry: An Introduction to General
Relativity},
Addison-Wesley (2004).

\bibitem{Rovelli}
C.~Rovelli,
\emph{Quantum Gravity},
Cambridge University Press (2004).

\bibitem{HR2012a}
S.~F.~Hassan and R.~A.~Rosen,
\emph{Bimetric gravity from ghost-free massive gravity},
JHEP \textbf{02} (2012) 126.

\bibitem{HR2012b}
S.~F.~Hassan and R.~A.~Rosen,
\emph{Confirmation of the secondary constraint and absence of
ghost in massive gravity and bimetric gravity},
JHEP \textbf{04} (2012) 123.

\bibitem{Alexandrov2013}
S.~Alexandrov,
\emph{Canonical structure of tetrad bimetric gravity},
Gen.\ Rel.\ Grav.\ \textbf{46} (2014) 1639,
arXiv:1308.6586.

\bibitem{LIGO2017}
B.~P.~Abbott et al.\ (LIGO Scientific and Virgo
Collaborations),
\emph{GW170817: Observation of Gravitational Waves from a
Binary Neutron Star Inspiral},
Phys.\ Rev.\ Lett.\ \textbf{119} (2017) 161101.

\bibitem{BD1972}
D.~G.~Boulware and S.~Deser,
\emph{Can gravitation have a finite range?}
Phys.\ Rev.\ D \textbf{6} (1972) 3368.

\bibitem{Nelson2026}
S.~Nelson,
\emph{Gravity Field Dynamics: Framework Essentials},
Zenodo (2026), DOI:~[to be inserted].

\bibitem{Nelson2026b}
S.~Nelson,
\emph{Gravity Field Dynamics and Bimetric Gravity: Dual
Tetrahedral Geometry as the Origin of the Hassan--Rosen
Interaction Parameters},
companion letter (2026).

\end{thebibliography}

\end{document}
